% ============================================================
% Presentación oral tentativa --- SOLABIMA 2026
% Biodiversidad en paisajes sonoros: detección de señales
% bioacústicas mediante aprendizaje automático no supervisado
%
% NOTA: presentación TENTATIVA. Autores/afiliaciones se dejaron
% como placeholder editable (ver \author y \institute abajo) --
% no se copiaron del resumen enviado a SOLABIMA a propósito.
% Grabadora corregida a Wildlife Acoustics Song Meter Micro 2
% (el resumen enviado decía "AudioMoth", dato desactualizado;
% la versión correcta está en el documento técnico/personal).
% ============================================================

\documentclass[11pt, aspectratio=169]{beamer}

\usepackage[utf8]{inputenc}
\usepackage[T1]{fontenc}
\usepackage[spanish, es-tabla]{babel}
\usepackage{amsmath, amssymb}
\usepackage{booktabs}
\usepackage{graphicx}
\usepackage{tikz}
\usetikzlibrary{arrows.meta, positioning}

% ---------- Tema: formal, sobrio, académico ----------
\usetheme{Madrid}
\usecolortheme{seahorse}
\setbeamertemplate{navigation symbols}{}
\setbeamertemplate{footline}[frame number]
\setbeamercovered{transparent}

% ---------- Datos de la charla (EDITAR) ----------
\title[Paisajes sonoros \& aprendizaje no supervisado]{%
  Biodiversidad en paisajes sonoros: detección de señales\\
  bioacústicas mediante aprendizaje automático no supervisado}
\subtitle{Prueba de concepto --- Reserva Natural Tapytá, Paraguay}

% TODO: completar autores/afiliaciones definitivos antes de enviar
\author[A.~Rojas et al.]{Alcides Rojas \and colaboradores}
\institute[]{Facultad Politécnica, Universidad Nacional de Asunción\\
  \textit{colaboración con Fundación Moisés Bertoni y otras instituciones}}
\date{SOLABIMA 2026}

\begin{document}

% ============================================================
\begin{frame}
  \titlepage
  % TODO: logos institucionales (subir archivos a una carpeta logos/ y descomentar).
  % ARASY no tiene logo propio (confirmado por el usuario y por conacyt.gov.py/arasy).
  % Sugeridos: CONACYT Paraguay (financia PINV01-530 y ARASY ESTR01-23),
  % FP-UNA/NIDTEC (institucion ejecutora), Fundacion Moises Bertoni, SOLABIMA.
  % \begin{center}
  %   \includegraphics[height=1cm]{logos/logo_conacyt.png}\hspace{1em}
  %   \includegraphics[height=1cm]{logos/logo_fpuna.png}\hspace{1em}
  %   \includegraphics[height=1cm]{logos/logo_fmb.png}
  % \end{center}
\end{frame}
% ============================================================

% ============================================================
\begin{frame}{Contexto y motivación}
% ============================================================
  \begin{itemize}
    \item El \textbf{monitoreo acústico pasivo (PAM)} genera volúmenes de audio
      que superan la capacidad de análisis manual o convencional.
    \item El paisaje sonoro (\textit{soundscape}) puede funcionar como
      \textbf{proxy de biodiversidad}: refleja la actividad de especies
      vocalmente activas sin necesidad de identificación experta previa.
    \item \textbf{Pregunta:} ¿puede el aprendizaje no supervisado detectar y
      caracterizar señales bioacústicas relevantes sin etiquetas previas?
  \end{itemize}
\end{frame}

% ============================================================
\begin{frame}{Objetivo}
% ============================================================
  \begin{block}{Objetivo general}
    Evaluar el potencial del aprendizaje no supervisado para detectar y
    caracterizar señales acústicas como proxy de biodiversidad, como prueba
    de concepto (PoC) en un subconjunto del monitoreo de la Reserva Natural
    Tapytá.
  \end{block}
  \vspace{0.5em}
  \begin{itemize}
    \item Identificar agrupamientos acústicos (clusters) sin supervisión.
    \item Comparar la complejidad acústica entre hábitats.
    \item Caracterizar patrones espaciotemporales de vocalización de anfibios.
  \end{itemize}
\end{frame}

% ============================================================
\begin{frame}{Área de estudio}
% ============================================================
  \begin{columns}[T]
    \begin{column}{0.55\textwidth}
      \textbf{Reserva Natural Tapytá}\\
      Bosque Atlántico del Alto Paraná, Paraguay
      \vspace{0.5em}
      \begin{itemize}
        \item 20 sitios en 3 hábitats:
          \begin{itemize}
            \item Bosque (BO, $n=10$)
            \item Eucaliptal (EU, $n=5$)
            \item Pastizal (PA, $n=5$)
          \end{itemize}
        \item Campaña: 11~nov -- 8~dic 2024 (temporada húmeda subtropical,
          inicio de vocalización de anfibios).
        \item Ventana horaria analizada: 18:00--06:00~h.
      \end{itemize}
    \end{column}
    \begin{column}{0.45\textwidth}
      \centering
      % TODO: reemplazar por mapa de sitios si se dispone de uno
      \textit{[Mapa de sitios --- pendiente]}
    \end{column}
  \end{columns}
\end{frame}

% ============================================================
\begin{frame}{Datos}
% ============================================================
  \begin{itemize}
    \item Grabadoras \textbf{Wildlife Acoustics Song Meter Micro~2}
      (WAV, 24{,}000~Hz, mono, 24~bits), activación horaria automática.
    \item $\approx$1{,}852{,}000 segmentos de 4~s (2~s de solapamiento)
      tras el filtro horario.
    \item Volumen total de referencia del estudio completo:
      $\sim$20~TB ($>$10 años de audio en tiempo real, múltiples sitios) ---
      esta PoC usa una submuestra (campaña nov--dic 2024).
  \end{itemize}
\end{frame}

% ============================================================
\begin{frame}{Metodología --- pipeline}
% ============================================================
  \centering
  \begin{tikzpicture}[
      node distance=6mm and 6mm,
      box/.style={rectangle, draw, rounded corners, align=center,
                  minimum height=1.1cm, text width=2.5cm, font=\scriptsize,
                  fill=blue!8},
      arr/.style={-{Stealth[length=2mm]}, thick}
    ]
    \node[box] (seg) {Segmentación\\4~s / 2~s solap.};
    \node[box, right=of seg] (aten) {Atenuación\\adaptativa\\insectos (6--12~kHz)};
    \node[box, right=of aten] (mfcc) {MFCC\\(20 coef.)};
    \node[box, right=of mfcc] (pca) {RobustScaler\\+ PCA (95\%)};
    \node[box, right=of pca] (umap) {UMAP\\$\mathbb{R}^d \to \mathbb{R}^3$};
    \node[box, right=of umap] (hdb) {HDBSCAN\\clustering};

    \draw[arr] (seg) -- (aten);
    \draw[arr] (aten) -- (mfcc);
    \draw[arr] (mfcc) -- (pca);
    \draw[arr] (pca) -- (umap);
    \draw[arr] (umap) -- (hdb);
  \end{tikzpicture}

  \vspace{1em}
  \begin{columns}[T]
    \begin{column}{0.48\textwidth}
      \small
      \textbf{UMAP}
      \begin{itemize}\footnotesize
        \item \texttt{n\_neighbors} = 60
        \item \texttt{min\_dist} = 0{,}3
        \item métrica: coseno
        \item 3 componentes
      \end{itemize}
    \end{column}
    \begin{column}{0.48\textwidth}
      \small
      \textbf{HDBSCAN}
      \begin{itemize}\footnotesize
        \item \texttt{min\_cluster\_size} = 800
        \item \texttt{min\_samples} = 65
        \item \texttt{epsilon} = 0{,}03
        \item método EOM
      \end{itemize}
    \end{column}
  \end{columns}
\end{frame}

% ============================================================
\begin{frame}{Refinamiento jerárquico}
% ============================================================
  \begin{itemize}
    \item Un cluster es candidato a subdivisión si:
      $n \geq 1{,}000$ segmentos, distancia media al centroide UMAP
      $\bar{d} \geq 0{,}8$, y percentil 90 $\geq 1{,}2$.
    \item \textbf{161 de 265 clusters (66{,}8\%)} fueron refinados,
      revelando subgrupos heterogéneos dentro de agrupamientos iniciales.
  \end{itemize}
  \begin{figure}
    \centering
    \includegraphics[width=0.75\textwidth]{figs/FigC_Subclustering_AntesDespues.jpg}
    \caption{\footnotesize Cluster 0 de PA-17 antes/después del subclustering.}
  \end{figure}
\end{frame}

% ============================================================
\begin{frame}{Verificación de consistencia \textnormal{\small(en revisión)}}
% ============================================================
  % TODO (pendiente esta semana): decidir si esto entra en la presentación.
  % Corresponde a Figura 5 (FigF1_Centroides_7subclusters, guia personal) y
  % Figura 8 (FigG1_Centroides_8clusters, guia personal) del documento tecnico/
  % personal: demasiados espectrogramas acumulados en un solo grafico y poca
  % diferencia visible entre ellos. Se deja el espacio de la figura vacio a
  % proposito hasta mejorar/simplificar el grafico o decidir si vale la pena.
  \begin{itemize}
    \item Los 7 sub-clusters del Cluster~0 muestran patrones espectrales
      heterogéneos (mezcla variable de vocalización de anfibio y coro de
      insectos).
    \item Los mismos tipos acústicos (vocalización de anfibio, coro de
      insectos) se reconocen de forma consistente en los tres hábitats
      (PA-17, EU-16, BO-31) --- evidencia de que el pipeline captura tipos
      sonoros reales, no artefactos de un sitio particular.
  \end{itemize}
  \vspace{1em}
  \begin{center}
    \fbox{\parbox{0.8\textwidth}{\centering\vspace{2em}
      \textit{[Figura pendiente de revisión --- ver nota de sesión]}
      \vspace{2em}}}
  \end{center}
  % \includegraphics[width=0.75\textwidth]{figs/FigF1_Centroides_7subclusters.jpg}
  % \includegraphics[width=0.75\textwidth]{figs/FigG1_Centroides_8clusters.jpg}
\end{frame}

% ============================================================
\begin{frame}{Resultados globales}
% ============================================================
  \begin{table}
    \centering
    \small
    \begin{tabular}{lcccc}
      \toprule
      \textbf{Hábitat} & \textbf{Clusters/sitio} & \textbf{\% ruido} & \textbf{Silueta} & \textbf{Davies--Bouldin} \\
      \midrule
      Pastizal   & 33{,}2 & 21{,}0\% & 0{,}437 & 0{,}810 \\
      Bosque     &  8{,}1 &  3{,}1\% & 0{,}368 & 0{,}643 \\
      Eucaliptal &  3{,}6 &  0{,}0\% & 0{,}404 & 0{,}505 \\
      \midrule
      \textbf{Global (265 clusters)} & \multicolumn{3}{l}{silueta 0{,}394 \quad ruido 6{,}8\%} & \\
      \bottomrule
    \end{tabular}
  \end{table}
  \begin{figure}
    \centering
    \includegraphics[width=0.7\textwidth]{figs/FigD_Violin_Comparacion_Habitats.jpg}
  \end{figure}
\end{frame}

% ============================================================
\begin{frame}{Diversidad acústica: marco metodológico}
% ============================================================
  Para cada sitio $s$ con $n_s$ segmentos y abundancias relativas
  $q_{k,s} = n_{k,s}/n_s$ por tipo acústico $k$:
  \begin{block}{Riqueza}
    $R_s = |K_s|$ \quad (número de tipos acústicos distintos --- valores ya
    mostrados en la tabla de la diapositiva anterior)
  \end{block}
  \begin{block}{Entropía (Shannon)}
    $\displaystyle H_s = -\sum_{k \in K_s} q_{k,s} \log q_{k,s}$
  \end{block}
  \vspace{0.3em}
  Las diferencias entre hábitats se evalúan mediante \textbf{pruebas de
  permutación} (10.000 réplicas) de $R_s$, $H_s$, $D_s$ (dominancia).

  \vspace{0.5em}
  {\footnotesize \textit{Nota: $R_s$ corresponde a los valores ya reportados;
  el cálculo de $H_s$/$D_s$ y el resultado formal de la prueba de permutación
  quedan pendientes (no están aún en el pipeline procesado).}\par}
\end{frame}

% ============================================================
\begin{frame}{Estructura del espacio acústico (UMAP)}
% ============================================================
  \begin{figure}
    \centering
    \includegraphics[width=0.95\textwidth]{figs/FigA_UMAP_Panel_3Habitats.jpg}
    \caption{\footnotesize Proyecciones UMAP 2D --- PA-17 (pastizal), EU-16
      (eucaliptal), BO-31 (bosque). Gris = ruido HDBSCAN.}
  \end{figure}
\end{frame}

% ============================================================
\begin{frame}{Caso de estudio: PA-17 (pastizal)}
% ============================================================
  \small
  \begin{table}
    \centering
    \footnotesize
    \begin{tabular}{clc}
      \toprule
      \textbf{Cluster} & \textbf{Categoría} & \textbf{$N$ segs.} \\
      \midrule
      Cl.\,1--3 & Vocalización de anfibio & 1{,}450 / 1{,}454 / 5{,}809 \\
      Cl.\,39   & Canto de ave & 3{,}528 \\
      Cl.\,17   & Coro de insectos & 10{,}884 \\
      Cl.\,38   & Lluvia / tormenta & 1{,}413 \\
      \bottomrule
    \end{tabular}
  \end{table}
  \begin{figure}
    \centering
    \includegraphics[width=0.85\textwidth]{figs/FigB_Panel_Espectrogramas_4Arquetipos.jpg}
  \end{figure}
  {\footnotesize Identificación por inspección espectral/auditiva con
  señuelos etiquetados, según marco de Thomas et al. (2022).\par}
\end{frame}

% ============================================================
\begin{frame}{Hallazgo principal: partición temporal del nicho acústico}
% ============================================================
  \begin{itemize}
    \item Tres clusters de vocalización de anfibio en PA-17,
      \textbf{espectralmente equivalentes} ($f_{\text{dom}} \approx 2{,}850$~Hz,
      BW $\approx$ 1{,}400--3{,}200~Hz).
    \item \textbf{Segregación temporal escalonada:} picos a las
      \textbf{18:00, 21:00 y 00:00~h}.
    \item Consistente con partición temporal del nicho acústico descrita
      en anuros del Bosque Atlántico brasileño (Donnelly et al., 2026).
  \end{itemize}
  \begin{figure}
    \centering
    \includegraphics[width=0.85\textwidth]{figs/FigJ_Actividad_Temporal_PA17.jpg}
  \end{figure}
\end{frame}

% ============================================================
\begin{frame}{Equivalencia espectral, divergencia temporal}
% ============================================================
  \begin{figure}
    \centering
    \includegraphics[width=0.8\textwidth]{figs/FigK_Comparacion_Espectral_Anfibios_PA17.jpg}
    \caption{\footnotesize Huellas espectrales de los tres clusters de
      anfibio: prácticamente idénticas. La diferencia está en el patrón
      temporal, no en el contenido espectral.}
  \end{figure}
\end{frame}

% ============================================================
\begin{frame}{Conclusiones}
% ============================================================
  \begin{enumerate}
    \item El pipeline no supervisado (MFCC + PCA + UMAP + HDBSCAN) produce
      clustering robusto en 20 sitios, sin ajuste manual por sitio.
    \item La complejidad acústica difiere marcadamente entre hábitats:
      pastizal $\gg$ bosque $\gg$ eucaliptal --- posible señal de
      empobrecimiento en monocultivo.
    \item Evidencia de \textbf{partición temporal del nicho acústico} entre
      clusters de anfibio espectralmente equivalentes.
    \item La similitud espectral no implica equivalencia temporal: el
      análisis combinado es necesario para caracterizar los grupos acústicos.
    \item El aprendizaje no supervisado muestra potencial claro para estudios
      tempranos de paisajes sonoros a gran escala.
  \end{enumerate}
\end{frame}

% ============================================================
\begin{frame}{Agradecimientos}
% ============================================================
  \begin{itemize}
    \item Financiamiento: \textbf{CONACYT Paraguay} --- Programa Nacional de
      Incentivo a la Investigación (\textbf{PINV01-530}) y proyecto
      estratégico \textbf{ARASY} (ESTR01-23): \textit{``Potenciamiento de las
      capacidades de modelamiento matemático computacional para la tierra y
      la vida''} (Programa PROCIENCIA~II / FEEI, ejecutado por FP-UNA/NIDTEC).
    \item Colaboración: Fundación Moisés Bertoni y equipo de campo en
      Reserva Natural Tapytá.
    \item Código disponible en:\\
      \url{https://github.com/AlcidesArielRojas/analisis_bioacustico_anfibios}
  \end{itemize}
  \vspace{1em}
  % TODO: logos institucionales (misma carpeta logos/ que la portada).
  % ARASY no tiene logo propio (confirmado) -> se representa con el logo de
  % CONACYT (financia PINV y ARASY) y el de FP-UNA/NIDTEC (institucion ejecutora).
  % \begin{center}
  %   \includegraphics[height=1cm]{logos/logo_conacyt.png}\hspace{1em}
  %   \includegraphics[height=1cm]{logos/logo_fpuna.png}\hspace{1em}
  %   \includegraphics[height=1cm]{logos/logo_fmb.png}\hspace{1em}
  %   \includegraphics[height=1cm]{logos/logo_solabima.png}
  % \end{center}
\end{frame}

% ============================================================
\begin{frame}
  \centering
  \Large \textbf{¡Gracias!}\\[1em]
  \normalsize Preguntas y comentarios
\end{frame}
% ============================================================

\end{document}
